../cictikz/src/cictikz/data/tex/cictikz_preamble.tex

% A single carrier link, end to end.
%
% One carrier, and the information rides on its amplitude and phase. The
% modulator turns those two into an I and a Q, the transmitter puts them
% on the carrier, and the receiver undoes both steps. Drawn beside the
% OFDM figure so the two can be compared: the blocks are the same, only
% what arrives at the left hand edge differs.

\begin{tikzpicture}[thick]
  %%%%%%%%%%%%%%%%%%%%%%%%%%%%%%%%%%%%%%%%%%%%%%%%%%%%%%%%%%%%%%%%%%%%%%
%% Shared pieces for the l04_afe signal flow graphs, l4_first_order and
%% l4_biquad.  The biquad is meant to read as the first-order graph with a
%% second integrator added, so the summing node and the 1/s block have to be
%% the same size and shape in both.  They are the same size in the original
%% artwork too.
%%
%% Names are prefixed and letters only, and computed coordinates are braced:
%% TikZ scans a coordinate for its closing parenthesis, so a '(' inside an
%% expression ends it early.
%%%%%%%%%%%%%%%%%%%%%%%%%%%%%%%%%%%%%%%%%%%%%%%%%%%%%%%%%%%%%%%%%%%%%%

\newcommand{\sfgR}{0.25}    % summing node radius
\newcommand{\sfgBw}{1.0}    % half width of the 1/s block
\newcommand{\sfgBh}{0.65}   % half height of the 1/s block

% A summing node at (#1,#2).
\newcommand{\sfgSum}[2]{%
  \draw (#1,#2) circle (\sfgR);
  \draw ({#1-0.17},#2) -- ({#1+0.17},#2);
  \draw (#1,{#2-0.17}) -- (#1,{#2+0.17});
}

% A block centred at (#1,#2) carrying #3.
\newcommand{\sfgBox}[3]{%
  \draw ({#1-\sfgBw},{#2-\sfgBh}) rectangle ({#1+\sfgBw},{#2+\sfgBh});
  \node at (#1,#2) {#3};
}

\newcommand{\sfgDot}[2]{\fill (#1,#2) circle (0.07);}


  \def\yt{0.95}
  \def\yb{-0.95}

  % ---- transmitter ----
  \draw[->] (-2.4,\yt) -- (0,\yt);
  \draw[->] (-2.4,\yb) -- (0,\yb);
  \node[armygreen, anchor=east] at (-2.5,\yt) {$A_m(t)$};
  \node[armygreen, anchor=east] at (-2.5,\yb) {$\phi_m(t)$};

  \sfgTall{0}{-1.5}{1.3}{1.5}{Mod}
  \draw[->] (1.3,\yt) -- (2.7,\yt) node[midway, above] {$I$};
  \draw[->] (1.3,\yb) -- (2.7,\yb) node[midway, above] {$Q$};
  \sfgTall{2.7}{-1.5}{4.0}{1.5}{TX}
  \draw (4.0,0) -- (5.0,0);
  \sfgAntenna{5.0}{0}

  % ---- receiver ----
  \sfgAntenna{7.4}{0}
  \draw (7.4,0) -- (8.4,0);
  \sfgTall{8.4}{-1.5}{9.7}{1.5}{RX}
  \draw[->] (9.7,\yt) -- (11.1,\yt) node[midway, above] {$I$};
  \draw[->] (9.7,\yb) -- (11.1,\yb) node[midway, above] {$Q$};
  \sfgTall{11.1}{-1.5}{12.4}{1.5}{De\\ mod}
  \draw[->] (12.4,\yt) -- (13.4,\yt);
  \draw[->] (12.4,\yb) -- (13.4,\yb);
  \node[armygreen, anchor=west] at (13.5,\yt) {$A_m(t)$};
  \node[armygreen, anchor=west] at (13.5,\yb) {$\phi_m(t)$};
\end{tikzpicture}
\end{document}
